\subsection{Degree preservation, closed-model competitors, and grouping}
The following lemmas make explicit the quantitative and algebraic
steps in the proof of Theorem~\ref{thm:joint}. In particular, regularity
of the competing channels is a conclusion of the small-error argument,
not an additional hypothesis on the competitor.

\begin{lemma}[The two leading-coefficient homotopies]\label{lem:leading-homotopies}
In the notation of \eqref{eq:polynomial-gauge}, put
\[
 h=\max_f\|h_f[r]\|_{\rm coef},\qquad
 e=\|H\|_{\rm coef},\qquad B=B_d(P).
\]
If $h\le1/4$ and $Be\le1/4$, the leading coefficients of
\[
 M+u(\overline M_r-M),\quad 0\le u\le1,
 \qquad\text{and}\qquad \overline M_r+sH,\quad0\le s\le1,
\]
are invertible. Their inverse norms are at most $2B$ and $4B$,
respectively. Each determinant has degree exactly $kd$.
In the observation comparison it suffices that
\begin{equation}\label{eq:quantitative-leading-smallness}
 C_0\delta/\tau(P)\le1/4,\qquad C_1B_d(P)\delta\le1/4,
\end{equation}
where $\|h_f[r]\|_{\rm coef}\le C_0\delta/\tau(P)$ and
$\|H\|_{\rm coef}\le C_1\delta$ are the interpolation bounds in that
proof. The constants involve only the fixed scalar experiment.
\end{lemma}
\begin{proof}
Write $h_{f,d}$ for the leading coefficient of $h_f[r]$. The inverse of
the leading matrix in the first homotopy is
\[
 \sum_f\frac{W_f}{1+u h_{f,d}}.
\]
All denominators have modulus at least $3/4$, so its norm is at most
$4B/3$. If $L$ is the leading matrix at $u=1$, then
$\|L^{-1}sH_d\|\le(4/3)Be\le1/3$. A Neumann series proves the
second assertion, with the stated bounds. The coefficient of $z^{kd}$
in the determinant of any $k\times k$ degree-$d$ matrix polynomial
is the determinant of its leading coefficient. It is nonzero
throughout both homotopies. Substitution gives
\eqref{eq:quantitative-leading-smallness}.
\end{proof}

\begin{lemma}[The reduced competitor is the target]\label{lem:competitor-rank}
Let $Q_1\in\R^{m_1\times k}$ and $Q_2\in\R^{m_2\times k}$ have
orthonormal columns, and let a competitor in $\cK_d$ have numerator
$K'=U'\diag(\alpha_b'f_b')V'^{\mathsf T}$. If the leading coefficient
of $M'=Q_1^{\mathsf T}K'Q_2$ is invertible, then both
$S'=Q_1^{\mathsf T}U'$ and $T'=Q_2^{\mathsf T}V'$ are invertible, and
\begin{equation}\label{eq:competitor-determinant-exact}
 \det M'(z)=\det S'\det T'\prod_{b=1}^k\alpha_b'f_b'(z).
\end{equation}
In particular the $kd$ roots counted by the homotopy argument are
exactly the competitor's component-root multiset, including all
multiplicities. Its original channels necessarily have full column
rank, even if this was not assumed when introducing the competitor.
\end{lemma}
\begin{proof}
The leading matrix is the product of the three square matrices
$S'$, $\diag(\alpha')$, and $T'^{\mathsf T}$. All weights are positive.
Invertibility of the product forces both outer factors to be invertible.
Taking determinants of the polynomial factorization proves
\eqref{eq:competitor-determinant-exact}. Since the $f_b'$ are monic of
degree $d$ and the scalar prefactor is nonzero, no root is cancelled
and no leading degree is lost. A rank-deficient $U'$ or $V'$ could not
have an invertible $k\times k$ compression.
\end{proof}

\begin{lemma}[Componentwise clusters]\label{lem:component-clusters}
For each scalar component polynomial $f$, partition its roots, counted
with multiplicity, into clusters of size at most $m$. Assume roots in
different clusters of the \emph{same} polynomial are separated by at
least $\nu>0$. No separation or common partition is imposed between
different component polynomials. On a fixed disk containing the roots,
\begin{equation}\label{eq:componentwise-scalar-lower}
 |f(z)|\ge c_\nu\min\{1,\dist(z,Z(f))^m\}
       \ge c_\nu\min\{1,\dist(z,\Lambda)^m\}.
\end{equation}
The constant is uniform when gaps inside a cluster vanish. Accordingly
the two exclusions in Theorem~\ref{thm:joint} use radius
$C_\nu(\|r\|_{\rm coef}+B_d\delta)^{1/m}$, with no cross-component
gap in the constant.
\end{lemma}
\begin{proof}
Put $a=\min\{1,\nu/4\}$. If the distance from $z$ to every root of
$f$ is at least $a$, then $|f(z)|\ge a^d$. Otherwise select a nearest
root and its cluster, containing $s\le m$ roots. Every root in another
cluster has distance at least $\nu-a\ge3\nu/4$ from $z$. Each of the
$s$ factors inside the chosen cluster has modulus at least
$\dist(z,Z(f))$. This gives
$|f(z)|\ge(3\nu/4)^{d-s}\dist(z,Z(f))^s$.
Since the last distance is less than one, its $s$th power is at least
its $m$th power. Taking the minimum of these finitely many constants
proves the first bound; $Z(f)\subset\Lambda$ gives the second.
The scalar gauge error is bounded by $C\|r\|_{\rm coef}$ and the
matrix error by $C\delta$. Applying \eqref{eq:gauge-resolved} with
this lower bound gives both homotopy exclusions. The leading-degree
lemma and the original contour matching argument then apply unchanged.
For simple roots take singleton clusters, so $m=1$. Disjoint root
brackets in Lemma~\ref{lem:local-fibre} are likewise required only
within each scalar polynomial, never between components.
\end{proof}

\begin{lemma}[Observable grouping semicontinuity]\label{lem:observable-grouping}
The lower semicontinuity and fixed-partition continuity of $B_d$ in
Theorem~\ref{thm:joint} follow from the recovered coefficient algebra
alone. No choice or limiting subsequence of latent representations is
needed.
\end{lemma}
\begin{proof}
Let $P_n\to P$ inside $\mathcal R_d$. Normalization recovery and
locally continuous orthonormal frames give $A_{0,n}\to A_0$ and
$C_{j,n}\to C_j$. Every joint tuple of the $C_{j,n}$ is bounded,
since each coordinate is an eigenvalue of a bounded matrix. Choose a
unit common eigenvector for any sequence of such tuples. Along a
subsequence both the vector and its tuple converge; the limit vector
is nonzero and is a common eigenvector for the $C_j$ with the limiting
tuple. Thus every nearby joint tuple lies near one of the finitely many
joint tuples at $P$. This argument is entirely within observable
matrix space.

Choose a fixed real linear combination $C(c)$ separating the distinct
limit tuples, and disjoint contours around its eigenvalues. The Riesz
projector of $C_n(c)$ on a contour converges to $E_f$. It equals the
sum of all joint projectors whose full tuples approach $f$, even if
some of their values under $c$ coincide. Consequently
\[
 \sum_{g\text{ approaching }f} A_{0,n}^{-1}E_{g,n}
                       \longrightarrow A_0^{-1}E_f=W_f.
\]
The triangle inequality, summed over limit tuples, gives
$B_d(P)\le\liminf_nB_d(P_n)$. With a fixed equality partition no
cluster splits or merges, and each summand converges separately.
Finally $\eta_d\le\tau$ handles limits with $\tau=0$; at rank-loss
limits with $\tau>0$, the resolved inverse at $T_1$ gives
$\sigma_k(P_n(T_1))^{-1}\le CB_d(P_n)$, forcing $B_d(P_n)\to\infty$.
This also proves the asserted upper semicontinuity of the zero
extension of $\eta_d$.
\end{proof}
